\section{実験計画法3；二要因モデル(Between)}
\label{lesson13}

\subsection{授業内容}
\subsubsection{科目の中でのこのコマの位置づけ}
\begin{tabular}[t]{ccccccc}
    記述統計[4] & 確率[2] & モデル[7] & 推測・推定[1] & 意思決定[0] & 実践[2]
\end{tabular}
\subsubsection{概要}
前回に続いて，二要因モデルを考える。回帰分析(一般線形モデル)の観点からは，重回帰モデルと同じである。一要因モデルと同じく誤差を算出することができるが，二つの要因だけでは説明できない大きさが残ることから，交互作用効果を考える必要があることを理解する。交互作用はどのようなパターンで見られるかを考えることと，重回帰分析で交互作用が言及されなかったことについて多重共線性の問題について理解する。
\subsubsection{コマ主題細目}
\begin{description}
\item[二要因のモデル]二要因の実験計画法モデルを考えるために，数値例を導入する。まず数値を可視化し，次に要因ごとに周辺化することで要因の効果をどのように書き出すことができるかを考える。回帰式に書き起こすことで，重回帰分析と同じ形式Formになっていることを確認する。
\begin{flushright}
    $\to$データは\citet{Yamada2004}のP.184,表7.6.1を用いる。計算式もP.184--189を参照すると良い。
\end{flushright}

\item[効果の算出と平方和の分解]回帰式を参照しながら，平方和を分解する計算を進める。二つの要因それぞれの大きさを各要因の水準ごとに考えることで計算する。要因ごとに区切るのは周辺化と呼ばれる。可視化したグラフを参考にしながら，計算箇所を辿りながら考えることが重要である。最後に二つの効果を全体平均に加え，水準の平均と合致「しない」ことを確認する。
\begin{flushright}
    $\to$ \citet{Yamada2004}のP.184--189を参照すると良い。
\end{flushright}

\item[交互作用]合計が合致しなかったのは，組み合わせによる効果が発生していることによる。このような効果のことを交互作用interactionという。交互作用の大きさを見積もることで平方和を完全に分解できたことになる。交互作用はその出現の仕方に様々なパターンがあることを理解する。また組み合わせ爆発という言葉があるように，三要因以上の実験計画を立てると，考慮すべき交互作用項がどんどん増えていくので，実質的に研究として使える実験計画は2，3要因の計画にとどめておいた方が良いことを理解する。
\begin{flushright}
    $\to$\citet{Yamada2004}のP.190--193.\citet{kosugi2019}のP.139--142
\end{flushright}

\item[多重共線性]実験計画は回帰モデルと同じことであり，二要因モデルに交互作用を考えるのであれば，説明変数が連続的である重回帰分析においても交互作用効果を検証するべきであったと気づくことができる。ただし連続変数の場合，複数の要因による組み合わせの効果があるということは説明変数同士が相関していることを意味する。説明変数の相関が高い重回帰分析モデルは，推定値が正しく得られない多重共線性とよばれる問題があることがある。
\begin{flushright}
    $\to$多重共線性の問題については，\citet{清水201705}のP.86--89。数学的には\citet{南風原200206}のP.244--247を参考にすると良い。計算機科学的には\citet{松浦201610}のP.103--104の見方も面白い。
\end{flushright}


\end{description}
\subsubsection{キーワード}
\begin{itemize}
\item 主効果
\item 交互作用
\item 多重共線性
\end{itemize}
\subsection{授業情報}
\paragraph{コマの展開方法}
講義
\paragraph{標準シラバスにおける位置づけ}
科目番号4；心理学研究法；2.データを用いた実証的な思考方法；C データの統計的記述\\
科目番号5；心理学統計法；1.心理学で用いられる統計手法；H 推測統計:代表値と散布度をめぐって(2)
\subsubsection{予習・復習課題}
\paragraph{予習}
実験計画と回帰分析が同じである，という原理を十分に理解し授業に臨む必要がある。またその観点から，複数の変数で行う回帰分析，すなわち重回帰分析についても復習しておく必要がある。
\paragraph{復習}
主効果と交互作用を算出するための手続きをもう一度，自分の手で展開しておくと良い。また前回に続いて，身の回りのデータや公開されているデータを使って，簡単に計算してみると良い。計算に当たっては統計環境Rのlm関数を使ってみることで，練習あるいは検算に役立てることができる。